%===============================================================================
% SEÇÃO 2 — Formulação Matemática do Problema de Controle
%===============================================================================

\section{Formulação Matemática do Problema de Controle}
\label{sec:formulacao}

% -----------------------------------------------------------------------
% SUBSEÇÃO 2.1 — MODELO DO PROCESSO
% O que escrever:
%   - Equação química da DRM com entalpia (eq. numerada).
%   - Vetor de estados x = [Tr, X_CH4, r_H2/CO]^T.
%   - Entradas u = [P_forno, F_CH4, F_CO2]^T.
%   - Forma genérica: dx/dt = f(x,u,d), y = g(x).
%   - Citar: \citep{oliveira2021,escribano2025,mokashi2024}
% -----------------------------------------------------------------------

% -----------------------------------------------------------------------
% SUBSEÇÃO 2.2 — PROBLEMA DE CONTROLE
% O que escrever:
%   - Determinar u(t) sujeito a:
%     Tr em [650, 850] graus C; X_CH4 > 0,90; 1 <= F_CO2/F_CH4 <= 1,5.
%   - Citar: \citep{luyben2014,thermodynamics2025}
% -----------------------------------------------------------------------

% -----------------------------------------------------------------------
% SUBSEÇÃO 2.3 — CRITÉRIO DE DESEMPENHO
% O que escrever:
%   - ISE: J = integral ||e(t)||_Q^2 + ||Delta_u||_R^2 dt (eq. numerada).
%   - Explicar Q e R.
% -----------------------------------------------------------------------
